%% ============================================================
%%  SECCIÓN 4 — Funciones de Operadores, Hadamard y Teorema BCH
%% ============================================================

\section{Funciones Analíticas de Operadores, Álgebra de Lie y Lemas BCH}
\label{sec:funciones_operadores}

\subsection{Definición de Funciones de Operadores Mediante Series de Potencias}

En mecánica cuántica avanzada, es sumamente frecuente encontrar expresiones donde un operador figura como argumento de una función trascendente, tal como la exponencial de un operador de evolución temporal $\eu^{-\iu\hat{H}t/\hbar}$ o de traslación espacial $\eu^{-\iu a\hat{p}/\hbar}$ \cite{dirac, sakurai}.

\begin{teoremabox}[Definición: Función Analítica de un Operador]
    Sea $f(z)$ una función analítica de variable compleja con desarrollo en serie de Taylor alrededor del origen con radio de convergencia $R$:
    \begin{equation}
        f(z) = \sum_{n=0}^\infty c_n z^n, \quad c_n = \frac{f^{(n)}(0)}{n!}.
    \end{equation}
    Para un operador lineal continuo $\hat{A}$ actuando sobre $\mathcal{H}$ cuyo radio espectral sea menor que $R$, la función del operador $f(\hat{A})$ se define formalmente por la serie de potencias operacionales:
    \begin{equation}
        f(\hat{A}) \equiv \sum_{n=0}^\infty c_n \hat{A}^n = c_0 \hat{I} + c_1 \hat{A} + c_2 \hat{A}^2 + \dots
        \label{eq:operator_function_def}
    \end{equation}
\end{teoremabox}

En particular, la \textbf{exponencial de un operador} es una función entera (convergente en norma para cualquier operador lineal acotado\footnote{En rigor del análisis funcional, observables cuánticos fundamentales como la posición $\hat{x}$, el momento $\hat{p}$ o el Hamiltoniano $\hat{H}$ son operadores \textit{no acotados} en $\mathcal{L}^2(\mathbb{R})$, de modo que sus series de potencias no convergen en norma sobre la totalidad del espacio. Para tales operadores autoadjuntos, la exponencial unitaria $\eu^{-\iu t\hat{A}/\hbar}$ (como $\eu^{-\iu a\hat{p}/\hbar}$ o $\eu^{-\iu t\hat{H}/\hbar}$) queda rigurosamente definida sobre todo el espacio de Hilbert $\mathcal{H}$ mediante el cálculo funcional espectral y el célebre \textbf{Teorema de Stone} sobre grupos uniparamétricos fuertemente continuos de operadores unitarios \cite{ballentine}. La expansión formal en serie de Taylor es válida en dominios densos de vectores analíticos. En este cuaderno mantendremos el tratamiento operacional estándar utilizado en la física teórica \cite{sakurai, zettili}.}):
\begin{equation}
    \boxed{\eu^{\hat{A}} \equiv \sum_{n=0}^\infty \frac{1}{n!} \hat{A}^n = \hat{I} + \hat{A} + \frac{1}{2!}\hat{A}^2 + \frac{1}{3!}\hat{A}^3 + \dots}
    \label{eq:operator_exponential}
\end{equation}

\begin{alertabox}[Propiedades Críticas de la Exponencial de Operadores]
    \begin{enumerate}
        \item Si $\hat{A}$ y $\hat{B}$ \textbf{conmutan} ($[\hat{A}, \hat{B}] = 0$), entonces se comportan como escalares y se cumple la propiedad exponencial usual:
        \begin{equation}
            \eu^{\hat{A}}\eu^{\hat{B}} = \eu^{\hat{A}+\hat{B}} = \eu^{\hat{B}}\eu^{\hat{A}}.
        \end{equation}
        En particular, $\eu^{\hat{A}}\eu^{-\hat{A}} = \eu^{\hat{A}-\hat{A}} = \eu^{\hat{0}} = \hat{I}$, lo que demuestra que:
        \begin{equation}
            (\eu^{\hat{A}})^{-1} = \eu^{-\hat{A}}.
        \end{equation}
        \item Si $\hat{A}$ y $\hat{B}$ \textbf{NO conmutan} ($[\hat{A}, \hat{B}] \neq 0$), entonces:
        \begin{equation}
            \boxed{\eu^{\hat{A}}\eu^{\hat{B}} \neq \eu^{\hat{A}+\hat{B}}}.
        \end{equation}
        El orden de los factores altera el resultado. Para calcular el producto en este caso general, es indispensable recurrir al Álgebra de Lie y al Teorema de Baker--Campbell--Hausdorff.
        \item Si $\hat{A}$ es hermítico ($\hat{A}\adj = \hat{A}$), el operador $\hat{U} = \eu^{\iu\hat{A}}$ es \textbf{unitario}:
        \begin{equation}
            \hat{U}\adj = \left(\sum_{n=0}^\infty \frac{\iu^n}{n!}\hat{A}^n\right)\adj = \sum_{n=0}^\infty \frac{(-\iu)^n}{n!}(\hat{A}\adj)^n = \sum_{n=0}^\infty \frac{(-\iu)^n}{n!}\hat{A}^n = \eu^{-\iu\hat{A}} = \hat{U}^{-1}.
        \end{equation}
    \end{enumerate}
\end{alertabox}

\subsection{El Adjunto de un Operador Función y la Conmutación Cuántica}
\label{subsec:adjunto_operador_funcion}

En el manejo algebraico de funciones de operadores, surgen dos interrogantes cruciales para la formulación cuántica: ¿cómo se calcula el conjugado hermítico de $F(\hat{A})$ y qué sucede cuando un operador $\hat{B}$ conmuta o no con $\hat{A}$?

\begin{teoremabox}[Propiedades del Adjunto de un Operador Función]
    Sea $F(z) = \sum_{n=0}^\infty c_n z^n$ una función analítica de variable compleja con coeficientes $c_n \in \mathbb{C}$, y sea $\hat{A}$ un operador lineal continuo. Entonces:
    \begin{enumerate}
        \item \textbf{Adjunto General:} El adjunto hermítico del operador función $F(\hat{A})$ está dado por:
        \begin{equation}
            \boxed{\left[ F(\hat{A}) \right]\adj = F^*(\hat{A}\adj) \equiv \sum_{n=0}^\infty c_n^* (\hat{A}\adj)^n}.
            \label{eq:adjunto_funcion_general}
        \end{equation}
        donde $F^*(z) \equiv \sum_{n=0}^\infty c_n^* z^n$ es la función con los coeficientes complejos conjugados.
        \item \textbf{Caso de Operador Hermítico:} Si $\hat{A}$ es autoadjunto ($\hat{A}\adj = \hat{A}$):
        \begin{equation}
            \left[ F(\hat{A}) \right]\adj = F^*(\hat{A}).
        \end{equation}
        En particular, si además los coeficientes $c_n$ son números reales ($F^*(z) = F(z)$), entonces $F(\hat{A})$ es un operador estrictamente \textbf{hermítico}:
        \begin{equation}
            \left[ F(\hat{A}) \right]\adj = F(\hat{A}).
        \end{equation}
        \item \textbf{Funciones Exponenciales Complejas:} Para cualquier parámetro escalar $\alpha \in \mathbb{C}$:
        \begin{equation}
            \boxed{\left( \eu^{\iu\alpha\hat{A}} \right)\adj = \eu^{-\iu\alpha^*\hat{A}\adj}}.
            \label{eq:adjunto_exponencial}
        \end{equation}
        Si $\alpha \in \mathbb{R}$ y $\hat{A}$ es hermítico ($\hat{A}\adj = \hat{A}$), entonces:
        \begin{equation}
            \left( \eu^{\iu\alpha\hat{A}} \right)\adj = \eu^{-\iu\alpha\hat{A}} = \left( \eu^{\iu\alpha\hat{A}} \right)^{-1},
        \end{equation}
        lo que demuestra rigurosamente que $\hat{U} \equiv \eu^{\iu\alpha\hat{A}}$ es un \textbf{operador unitario}.
    \end{enumerate}
\end{teoremabox}

\subsubsection{No Conmutatividad con Funciones de Operadores}

Si dos observables conmutan ($[\hat{A}, \hat{B}] = 0$), entonces $\hat{B}$ conmuta con cualquier potencia entera $\hat{A}^n$, y por ende conmuta con cualquier función analítica:
\begin{equation}
    [\hat{A}, \hat{B}] = 0 \implies [\hat{B}, F(\hat{A})] = 0.
\end{equation}
Por el contrario, si los operadores \textbf{no conmutan} ($[\hat{A}, \hat{B}] \neq 0$), entonces en general $\hat{B}$ \textbf{no conmuta con la función del operador}:
\begin{equation}
    [\hat{B}, F(\hat{A})] \neq 0 \quad (\text{en general}).
\end{equation}
(Cabe advertir que la implicación estricta $[\hat{A}, \hat{B}] \neq 0 \implies [\hat{B}, F(\hat{A})] \neq 0$ no es universal: por ejemplo, tomando las matrices de Pauli con $\hat{A} = \sigma_x$, $\hat{B} = \sigma_z$ y la función $F(z) = z^2$, se tiene $[\sigma_x, \sigma_z] = -2\iu\sigma_y \neq 0$, pero $F(\sigma_x) = \sigma_x^2 = \hat{I}$ conmuta idénticamente con $\sigma_z$).
Para cuantificar este conmutador de forma analítica, observemos la acción de $\hat{B}$ sobre una potencia monomial $\hat{A}^n$:
\begin{equation}
    [\hat{B}, \hat{A}^n] = \sum_{k=0}^{n-1} \hat{A}^k [\hat{B}, \hat{A}] \hat{A}^{n-1-k}.
\end{equation}
En el caso físicamente fundamental en que el conmutador básico conmuta con $\hat{A}$ (es decir, $[\hat{A}, [\hat{A}, \hat{B}]] = 0$, como ocurre para el par canónico de posición y momento $[\hat{x}, \hat{p}_x] = \iu\hbar\hat{I}$), el término $[\hat{B}, \hat{A}]$ sale como constante de la sumatoria:
\begin{equation}
    [\hat{B}, \hat{A}^n] = [\hat{B}, \hat{A}] \sum_{k=0}^{n-1} \hat{A}^{n-1} = n\hat{A}^{n-1}[\hat{B}, \hat{A}] = -n\hat{A}^{n-1}[\hat{A}, \hat{B}].
\end{equation}
Multiplicando por los coeficientes de Taylor $c_n$ y sumando sobre $n$:
\begin{equation}
    [\hat{B}, F(\hat{A})] = \sum_{n=1}^\infty c_n [\hat{B}, \hat{A}^n] = -[\hat{A}, \hat{B}] \sum_{n=1}^\infty n c_n \hat{A}^{n-1} = \boxed{-[\hat{A}, \hat{B}]\,F'(\hat{A})},
    \label{eq:comm_b_f_a}
\end{equation}
donde $F'(z) = \frac{\dif F}{\dif z}$ es la derivada compleja ordinaria de la función. Esta elegante identidad operativa permite calcular de un solo vistazo conmutadores que de otra manera requerirían cálculos matriciales interminables.

\subsection{El Lema de Hadamard (Fórmula de Trenzado de Operadores)}

Consideremos la transformación de un operador $\hat{B}$ mediante una conjugación unitaria o transformación de semejanza exponencial generada por $\hat{A}$:
\begin{equation}
    \hat{F}(\lambda) \equiv \eu^{\lambda\hat{A}}\hat{B}\eu^{-\lambda\hat{A}},
    \label{eq:hadamard_ansatz}
\end{equation}
donde $\lambda$ es un parámetro escalar continuo (real o complejo, independiente de las coordenadas espaciales).

\begin{teoremabox}[Lema de Hadamard]
    Para dos operadores lineales cualesquiera $\hat{A}$ y $\hat{B}$:
    \begin{align}
        \eu^{\lambda\hat{A}}\hat{B}\eu^{-\lambda\hat{A}} &= \hat{B} + \lambda[\hat{A}, \hat{B}] + \frac{\lambda^2}{2!}[\hat{A}, [\hat{A}, \hat{B}]] + \frac{\lambda^3}{3!}[\hat{A}, [\hat{A}, [\hat{A}, \hat{B}]]] + \dots \nonumber\\
        &= \sum_{n=0}^\infty \frac{\lambda^n}{n!} \mathrm{ad}_{\hat{A}}^n(\hat{B}),
        \label{eq:hadamard_lemma}
    \end{align}
    donde el operador de la representación adjunta $\mathrm{ad}_{\hat{A}}$ se define recursivamente por:
    \begin{equation}
        \mathrm{ad}_{\hat{A}}^0(\hat{B}) = \hat{B}, \quad \mathrm{ad}_{\hat{A}}^1(\hat{B}) = [\hat{A}, \hat{B}], \quad \mathrm{ad}_{\hat{A}}^n(\hat{B}) = [\hat{A}, \mathrm{ad}_{\hat{A}}^{n-1}(\hat{B})].
    \end{equation}
    En particular, evaluando en el valor unitario del parámetro $\lambda=1$:
    \begin{equation}
        \boxed{\eu^{\hat{A}}\hat{B}\eu^{-\hat{A}} = \hat{B} + [\hat{A}, \hat{B}] + \frac{1}{2!}[\hat{A}, [\hat{A}, \hat{B}]] + \frac{1}{3!}[\hat{A}, [\hat{A}, [\hat{A}, \hat{B}]]] + \dots}
    \end{equation}
\end{teoremabox}

\textbf{Demostración Rigurosa (Método de la Función Generatriz Paramétrica):}\\
Puesto que $\hat{F}(\lambda)$ depende analíticamente del escalar $\lambda$, podemos expandirla en una serie de Maclaurin ordinaria alrededor de $\lambda = 0$:
\begin{equation}
    \hat{F}(\lambda) = \hat{F}(0) + \lambda\left.\frac{\dif\hat{F}}{\dif \lambda}\right|_{\lambda=0} + \frac{\lambda^2}{2!}\left.\frac{\dif^2\hat{F}}{\dif \lambda^2}\right|_{\lambda=0} + \dots = \sum_{n=0}^\infty \frac{\lambda^n}{n!} \left.\frac{\dif^n\hat{F}}{\dif \lambda^n}\right|_{\lambda=0}.
    \label{eq:hadamard_maclaurin}
\end{equation}
Evaluemos el término de orden cero en $\lambda = 0$:
\begin{equation}
    \hat{F}(0) = \eu^{0\hat{A}}\hat{B}\eu^{-0\hat{A}} = \hat{I}\hat{B}\hat{I} = \hat{B}.
\end{equation}
Calculemos ahora la primera derivada de $\hat{F}(\lambda)$ respecto a $\lambda$ empleando la regla del producto de Leibniz para operadores:
\begin{align}
    \frac{\dif\hat{F}}{\dif \lambda} &= \frac{\dif}{\dif \lambda}\left( \eu^{\lambda\hat{A}}\hat{B}\eu^{-\lambda\hat{A}} \right) \nonumber\\
    &= \left(\frac{\dif}{\dif \lambda}\eu^{\lambda\hat{A}}\right)\hat{B}\eu^{-\lambda\hat{A}} + \eu^{\lambda\hat{A}}\hat{B}\left(\frac{\dif}{\dif \lambda}\eu^{-\lambda\hat{A}}\right).
\end{align}
Derivando término a término la serie de la exponencial escalar-operador:
\begin{equation}
    \frac{\dif}{\dif \lambda}\eu^{\lambda\hat{A}} = \frac{\dif}{\dif \lambda}\sum_{n=0}^\infty \frac{\lambda^n}{n!}\hat{A}^n = \sum_{n=1}^\infty \frac{n \lambda^{n-1}}{n!}\hat{A}^n = \hat{A}\sum_{m=0}^\infty \frac{\lambda^m}{m!}\hat{A}^m = \hat{A}\eu^{\lambda\hat{A}} = \eu^{\lambda\hat{A}}\hat{A}.
\end{equation}
Análogamente, para la exponencial negativa:
\begin{equation}
    \frac{\dif}{\dif \lambda}\eu^{-\lambda\hat{A}} = -\hat{A}\eu^{-\lambda\hat{A}} = -\eu^{-\lambda\hat{A}}\hat{A}.
\end{equation}
Sustituyendo estas derivadas:
\begin{align}
    \frac{\dif\hat{F}}{\dif \lambda} &= \hat{A}\left(\eu^{\lambda\hat{A}}\hat{B}\eu^{-\lambda\hat{A}}\right) - \left(\eu^{\lambda\hat{A}}\hat{B}\eu^{-\lambda\hat{A}}\right)\hat{A} \nonumber\\
    &= \hat{A}\hat{F}(\lambda) - \hat{F}(\lambda)\hat{A} = [\hat{A}, \hat{F}(\lambda)].
    \label{eq:hadamard_first_diff}
\end{align}
Evaluando la primera derivada en $\lambda = 0$:
\begin{equation}
    \left.\frac{\dif\hat{F}}{\dif \lambda}\right|_{\lambda=0} = [\hat{A}, \hat{F}(0)] = [\hat{A}, \hat{B}].
\end{equation}
Derivemos nuevamente respecto a $\lambda$ para obtener el término de segundo orden:
\begin{equation}
    \frac{\dif^2\hat{F}}{\dif \lambda^2} = \frac{\dif}{\dif \lambda}\left( \frac{\dif\hat{F}}{\dif \lambda} \right) = \frac{\dif}{\dif \lambda}[\hat{A}, \hat{F}(\lambda)] = \left[\hat{A}, \frac{\dif\hat{F}}{\dif \lambda}\right].
\end{equation}
Evaluando en $\lambda = 0$:
\begin{equation}
    \left.\frac{\dif^2\hat{F}}{\dif \lambda^2}\right|_{\lambda=0} = \left[\hat{A}, \left.\frac{\dif\hat{F}}{\dif \lambda}\right|_{\lambda=0}\right] = [\hat{A}, [\hat{A}, \hat{B}]].
\end{equation}
Por inducción matemática inmediata, la $n$-ésima derivada evaluada en el origen es:
\begin{equation}
    \left.\frac{\dif^n\hat{F}}{\dif \lambda^n}\right|_{\lambda=0} = \mathrm{ad}_{\hat{A}}^n(\hat{B}) = \underbrace{[\hat{A}, [\hat{A}, \dots, [\hat{A}, \hat{B}]] \dots ]}_{n \text{ conmutadores anidados}}.
\end{equation}
Sustituyendo estos coeficientes en la serie de Taylor \eqref{eq:hadamard_maclaurin}, queda demostrado formalmente el Lema de Hadamard \cite{zettili}.

\subsection{El Teorema de Baker--Campbell--Hausdorff (BCH) Simplificado}

¿Cómo podemos calcular el producto de dos exponenciales $\eu^{\hat{A}}\eu^{\hat{B}}$ cuando los operadores no conmutan? El teorema general de Baker, Campbell y Hausdorff establece que existe un operador $\hat{C}$ tal que $\eu^{\hat{A}}\eu^{\hat{B}} = \eu^{\hat{C}}$, donde $\hat{C} = \hat{A} + \hat{B} + \frac{1}{2}[\hat{A}, \hat{B}] + \frac{1}{12}[\hat{A}, [\hat{A}, \hat{B}]] - \frac{1}{12}[\hat{B}, [\hat{A}, \hat{B}]] + \dots$

En física cuántica, existe un caso extraordinariamente frecuente y de vital importancia (llamado el \textbf{caso c-número}): aquel donde el conmutador $[\hat{A}, \hat{B}]$ conmuta tanto con $\hat{A}$ como con $\hat{B}$, es decir:
\begin{equation}
    [[\hat{A}, \hat{B}], \hat{A}] = 0 \quad \text{y} \quad [[\hat{A}, \hat{B}], \hat{B}] = 0.
    \label{eq:bch_hypothesis}
\end{equation}
Esto ocurre siempre que $[\hat{A}, \hat{B}] = c\hat{I}$ es un escalar (por ejemplo, el conmutador canónico de posición y momento $[\hat{x}, \hat{p}_x] = \iu\hbar\hat{I}$, o los operadores de creación y aniquilación $[\hat{a}, \hat{a}\adj] = \hat{I}$).

\begin{teoremabox}[Teorema de Baker--Campbell--Hausdorff Simplificado (Fórmula de Weyl)]
    Si dos operadores lineales $\hat{A}$ y $\hat{B}$ satisfacen la condición \eqref{eq:bch_hypothesis}, entonces:
    \begin{equation}
        \boxed{\eu^{\hat{A}}\eu^{\hat{B}} = \eu^{\hat{A}+\hat{B}}\eu^{\frac{1}{2}[\hat{A}, \hat{B}]}}.
        \label{eq:bch_simplified}
    \end{equation}
\end{teoremabox}

\textbf{Demostración Paso a Paso:}\\
Para demostrar esta relación, introducimos un operador dependiente de un parámetro continuo $t \in [0, 1]$:
\begin{equation}
    \hat{G}(t) \equiv \eu^{t\hat{A}}\eu^{t\hat{B}}.
    \label{eq:g_ansatz}
\end{equation}
Notemos que para $t = 0$: $\hat{G}(0) = \eu^{\hat{0}}\eu^{\hat{0}} = \hat{I}$, y para $t = 1$: $\hat{G}(1) = \eu^{\hat{A}}\eu^{\hat{B}}$ (que es la magnitud deseada).

Derivemos $\hat{G}(t)$ respecto al parámetro escalar $t$:
\begin{equation}
    \frac{\dif\hat{G}}{\dif t} = \left(\frac{\dif}{\dif t}\eu^{t\hat{A}}\right)\eu^{t\hat{B}} + \eu^{t\hat{A}}\left(\frac{\dif}{\dif t}\eu^{t\hat{B}}\right) = \hat{A}\eu^{t\hat{A}}\eu^{t\hat{B}} + \eu^{t\hat{A}}\hat{B}\eu^{t\hat{B}}.
\end{equation}

\begin{alertabox}[¿Por qué insertamos la identidad en la forma $\eu^{-t\hat{A}}\eu^{t\hat{A}}$?]
    En el primer término, $\eu^{t\hat{A}}\eu^{t\hat{B}} = \hat{G}(t)$, de modo que queda $\hat{A}\hat{G}(t)$.\\
    Sin embargo, en el segundo término, el operador $\hat{B}$ se encuentra \enquote{atrapado} en medio de las exponenciales: $\eu^{t\hat{A}}\hat{B}\eu^{t\hat{B}}$.\\
    \textbf{¡NO podemos simplemente sacar $\hat{B}$ hacia la izquierda ni hacia la derecha porque $\hat{B}$ no conmuta con $\eu^{t\hat{A}}$!}
    
    Para poder factorizar $\hat{G}(t) = \eu^{t\hat{A}}\eu^{t\hat{B}}$ por la derecha de toda la ecuación, insertamos estratégicamente el operador identidad $\hat{I} = \eu^{-t\hat{A}}\eu^{t\hat{A}}$ justo entre $\hat{B}$ y $\eu^{t\hat{B}}$:
    \begin{equation}
        \eu^{t\hat{A}}\hat{B}\eu^{t\hat{B}} = \eu^{t\hat{A}}\hat{B}\underbrace{(\eu^{-t\hat{A}}\eu^{t\hat{A}})}_{\hat{I}}\eu^{t\hat{B}} = \left(\eu^{t\hat{A}}\hat{B}\eu^{-t\hat{A}}\right)\left(\eu^{t\hat{A}}\eu^{t\hat{B}}\right) = \left(\eu^{t\hat{A}}\hat{B}\eu^{-t\hat{A}}\right)\hat{G}(t).
    \end{equation}
    ¡Observemos la genialidad de este paso algebraico! La expresión entre paréntesis es exactamente la estructura del \textbf{Lema de Hadamard} con parámetro $t$.
\end{alertabox}

Aplicando el Lema de Hadamard \eqref{eq:hadamard_lemma} a la cantidad $\eu^{t\hat{A}}\hat{B}\eu^{-t\hat{A}}$:
\begin{equation}
    \eu^{t\hat{A}}\hat{B}\eu^{-t\hat{A}} = \hat{B} + t[\hat{A}, \hat{B}] + \frac{t^2}{2!}[\hat{A}, [\hat{A}, \hat{B}]] + \dots
\end{equation}
Pero según nuestra hipótesis fundamental de trabajo \eqref{eq:bch_hypothesis}, el segundo conmutador y todos los superiores se anulan idénticamente: $[\hat{A}, [\hat{A}, \hat{B}]] = 0$. Por lo tanto, la serie infinita \textbf{se trunca exactamente en el segundo término}:
\begin{equation}
    \eu^{t\hat{A}}\hat{B}\eu^{-t\hat{A}} = \hat{B} + t[\hat{A}, \hat{B}].
\end{equation}
Sustituyendo esto en la derivada de $\hat{G}(t)$:
\begin{align}
    \frac{\dif\hat{G}}{\dif t} &= \hat{A}\hat{G}(t) + (\hat{B} + t[\hat{A}, \hat{B}])\hat{G}(t) \nonumber\\
    &= \left( \hat{A} + \hat{B} + t[\hat{A}, \hat{B}] \right)\hat{G}(t).
    \label{eq:diff_eq_g}
\end{align}
Denotemos al operador entre paréntesis por $\hat{M}(t) \equiv \hat{A} + \hat{B} + t[\hat{A}, \hat{B}]$. Como $[\hat{A}+\hat{B}, [\hat{A}, \hat{B}]] = 0$, el operador $\hat{M}(t)$ conmuta consigo mismo a diferentes instantes $t$ y $t'$:
\begin{equation}
    [\hat{M}(t), \hat{M}(t')] = [\hat{A}+\hat{B} + t[\hat{A},\hat{B}], \, \hat{A}+\hat{B} + t'[\hat{A},\hat{B}]] = 0.
\end{equation}
Debido a esta conmutación, la ecuación diferencial de operadores \eqref{eq:diff_eq_g} se integra exactamente como una ecuación diferencial escalar por separación de variables:
\begin{equation}
    \hat{G}(t) = \exp\left( \int_0^t \hat{M}(t')\,\dif t' \right)\hat{G}(0).
\end{equation}
Evaluemos la integral en el argumento de la exponencial:
\begin{equation}
    \int_0^t \left( \hat{A} + \hat{B} + t'[\hat{A}, \hat{B}] \right)\,\dif t' = t(\hat{A} + \hat{B}) + \frac{t^2}{2}[\hat{A}, \hat{B}].
\end{equation}
Puesto que $\hat{G}(0) = \hat{I}$:
\begin{equation}
    \hat{G}(t) = \exp\left( t(\hat{A} + \hat{B}) + \frac{1}{2}t^2[\hat{A}, \hat{B}] \right).
\end{equation}
Nuevamente, como el operador $[\hat{A}, \hat{B}]$ conmuta con $(\hat{A}+\hat{B})$, podemos separar la exponencial de la suma en el producto de exponenciales:
\begin{equation}
    \hat{G}(t) = \eu^{t(\hat{A}+\hat{B})}\eu^{\frac{1}{2}t^2[\hat{A}, \hat{B}]}.
\end{equation}
Evaluando finalmente en $t = 1$:
\begin{equation}
    \hat{G}(1) = \eu^{\hat{A}}\eu^{\hat{B}} = \eu^{\hat{A}+\hat{B}}\eu^{\frac{1}{2}[\hat{A}, \hat{B}]}.
\end{equation}
La demostración queda enteramente completada y justificada.

\subsection{BCH Simplificado (c-Número) frente a la Serie BCH General de Momento Angular}
\label{subsec:bch_simplificado_vs_general}

Un error pedagógico extraordinariamente común en los primeros cursos de física cuántica es asumir que la fórmula:
\begin{equation}
    \eu^{\hat{A}}\eu^{\hat{B}} = \eu^{\hat{A}+\hat{B} + \frac{1}{2}[\hat{A}, \hat{B}]}
    \label{eq:bch_simplificada_recap}
\end{equation}
es válida universalmente para cualquier par de operadores cuánticos. \textbf{Esto es falso}: la Ec.~\eqref{eq:bch_simplificada_recap} es un caso particular que depende de una condición sumamente restrictiva.

\begin{alertabox}[¿Cuándo es válida la fórmula simplificada de Baker--Campbell--Hausdorff?]
    La fórmula cerrada \eqref{eq:bch_simplificada_recap} exige estrictamente que el conmutador $[\hat{A}, \hat{B}]$ sea un \textbf{c-número} (un múltiplo escalar de la identidad $\alpha\hat{I}$), o al menos que conmute con ambos operadores:
    \begin{equation}
        [\hat{A}, [\hat{A}, \hat{B}]] = 0 \quad \text{y} \quad [\hat{B}, [\hat{A}, \hat{B}]] = 0.
    \end{equation}
    \begin{itemize}
        \item \textbf{El caso canónico de posición y momento $(\hat{x}, \hat{p}_x)$:}
        Aquí $[\hat{x}, \hat{p}_x] = \iu\hbar\hat{I}$. Dado que $\iu\hbar\hat{I}$ es proporcional al operador identidad, conmuta absolutamente con cualquier operador del universo:
        \begin{equation}
            [\hat{x}, \, \iu\hbar\hat{I}] = 0, \qquad [\hat{p}_x, \, \iu\hbar\hat{I}] = 0.
        \end{equation}
        Por lo tanto, todos los conmutadores anidados de segundo orden en adelante se anulan idénticamente, y la fórmula simplificada \eqref{eq:bch_simplificada_recap} es \textbf{rigurosamente exacta}.
        \item \textbf{El caso de momento angular $(\hat{L}_x, \hat{L}_y)$:}
        Consideremos ahora dos generadores de rotaciones espaciales. Su conmutador básico es el álgebra de Lie $\mathfrak{so}(3)$:
        \begin{equation}
            [\hat{L}_x, \hat{L}_y] = \iu\hbar\hat{L}_z.
        \end{equation}
        ¡El resultado no es un número; \textbf{es otro operador cuántico hermítico ($\hat{L}_z$)}!\\
        Evaluemos el siguiente conmutador anidado de la serie:
        \begin{equation}
            [\hat{L}_x, [\hat{L}_x, \hat{L}_y]] = [\hat{L}_x, \, \iu\hbar\hat{L}_z] = \iu\hbar[\hat{L}_x, \hat{L}_z] = \iu\hbar(-\iu\hbar\hat{L}_y) = +\hbar^2\hat{L}_y \neq 0.
        \end{equation}
        Análogamente, para el segundo operador:
        \begin{equation}
            [\hat{L}_y, [\hat{L}_x, \hat{L}_y]] = [\hat{L}_y, \, \iu\hbar\hat{L}_z] = \iu\hbar[\hat{L}_y, \hat{L}_z] = \iu\hbar(\iu\hbar\hat{L}_x) = -\hbar^2\hat{L}_x \neq 0.
        \end{equation}
        ¡Los conmutadores anidados nunca se anulan; se reproducen infinitamente rotando entre $\hat{L}_x$, $\hat{L}_y$ y $\hat{L}_z$!\\
        Por consiguiente, \textbf{la fórmula exponencial simplificada falla por completo para el momento angular}:
        \begin{equation}
            \eu^{-\frac{\iu}{\hbar}\alpha\hat{L}_x}\eu^{-\frac{\iu}{\hbar}\beta\hat{L}_y} \neq \eu^{-\frac{\iu}{\hbar}(\alpha\hat{L}_x + \beta\hat{L}_y) + \frac{1}{2}\left[-\frac{\iu}{\hbar}\alpha\hat{L}_x, -\frac{\iu}{\hbar}\beta\hat{L}_y\right]}.
            \label{eq:bch_falla_L}
        \end{equation}
    \end{itemize}
\end{alertabox}

\begin{teoremabox}[La Serie General de Baker--Campbell--Hausdorff (BCH)]
    Para dos operadores genéricos $\hat{A}$ y $\hat{B}$ que no satisfacen la hipótesis de truncamiento a primer orden, el logaritmo del producto exponencial viene dado por la serie infinita completa de conmutadores anidados:
    \begin{equation}
        \ln\left( \eu^{\hat{A}}\eu^{\hat{B}} \right) = \hat{A} + \hat{B} + \frac{1}{2}[\hat{A}, \hat{B}] + \frac{1}{12}[\hat{A}, [\hat{A}, \hat{B}]] - \frac{1}{12}[\hat{B}, [\hat{A}, \hat{B}]] - \frac{1}{24}[\hat{B}, [\hat{A}, [\hat{A}, \hat{B}]]] + \dots
        \label{eq:bch_general_series}
    \end{equation}
    En la teoría de grupos de Lie, esta serie describe la composición de dos rotaciones espaciales en el grupo $\mathrm{SO}(3)$: una rotación de ángulo $\alpha$ alrededor de $x$ seguida de una rotación $\beta$ alrededor de $y$ no es una simple suma vectorial de rotaciones, sino una nueva rotación alrededor de un eje oblicuo resultante de la no conmutatividad geométrica de las rotaciones en $\mathbb{R}^3$.
\end{teoremabox}

\begin{figure}[H]
\centering
\begin{tikzpicture}[
    x={(-0.50cm,-0.36cm)}, y={(1cm,0cm)}, z={(0cm,1cm)},
    scale=1.6, >=Stealth, font=\small]

  % ---------- Panel (a): primero R_x, luego R_y ----------
  \begin{scope}
    \draw[grisUNI!40] (0,0,0)--(-1.2,0,0);
    \draw[grisUNI!40] (0,0,0)--(0,-1.6,0);
    \draw[grisUNI!40] (0,0,0)--(0,0,-1.5);
    \draw[->,grisUNI] (0,0,0)--(1.8,0,0) node[below left]{$x$};
    \draw[->,grisUNI] (0,0,0)--(0,1.8,0) node[right]{$y$};
    \draw[->,grisUNI] (0,0,0)--(0,0,1.8) node[above]{$z$};
    \draw[->,thick,grisUNI!70!black] (0,0,0)--(1,0,1)
          node[above left]{$\vec v=(1,0,1)$};
    \draw[->,thick,azulUNI,dashed] (0,0,0)--(1,-1,0)
          node[left]{$(1,-1,0)$};
    \draw[->,very thick,rojoUNI] (0,0,0)--(0,-1,-1)
          node[below]{$(0,-1,-1)$};
    \node at (0,0.4,2.3)
          {\textbf{(a)} $\hat R_y(\tfrac{\pi}{2})\,\hat R_x(\tfrac{\pi}{2})\,\vec v$};
  \end{scope}

  % ---------- Panel (b): primero R_y, luego R_x ----------
  \begin{scope}[shift={(canvas cs:x=0cm,y=0cm)}]
    \pgftransformxshift{4.9cm}
    \draw[grisUNI!40] (0,0,0)--(-1.2,0,0);
    \draw[grisUNI!40] (0,0,0)--(0,-1.6,0);
    \draw[grisUNI!40] (0,0,0)--(0,0,-1.5);
    \draw[->,grisUNI] (0,0,0)--(1.8,0,0) node[below left]{$x$};
    \draw[->,grisUNI] (0,0,0)--(0,1.8,0) node[right]{$y$};
    \draw[->,grisUNI] (0,0,0)--(0,0,1.8) node[above]{$z$};
    \draw[->,thick,grisUNI!70!black] (0,0,0)--(1,0,1)
          node[above left]{$\vec v=(1,0,1)$};
    \draw[->,thick,azulUNI,dashed] (0,0,0)--(1,0,-1)
          node[below left]{$(1,0,-1)$};
    \draw[->,very thick,rojoUNI] (0,0,0)--(1,1,0)
          node[below right]{$(1,1,0)$};
    \node at (0,0.4,2.3)
          {\textbf{(b)} $\hat R_x(\tfrac{\pi}{2})\,\hat R_y(\tfrac{\pi}{2})\,\vec v$};
  \end{scope}
\end{tikzpicture}
\caption{Las rotaciones finitas no conmutan. Un mismo vector $\vec v=(1,0,1)$,
girado $\pi/2$ alrededor de $x$ y de $y$ en órdenes distintos, termina en
posiciones distintas (rojo). Es la versión geométrica de
$\eu^{-\iu\alpha\hat L_x/\hbar}\eu^{-\iu\beta\hat L_y/\hbar}\neq
\eu^{-\iu\beta\hat L_y/\hbar}\eu^{-\iu\alpha\hat L_x/\hbar}$ (ec.~\eqref{eq:bch_falla_L}); línea
discontinua: posición intermedia.}
\label{fig:rot-no-conmutan}
\end{figure}


\subsection{Ejemplo Resuelto 3: El Operador de Traslación Espacial y Serie de Taylor}

\begin{ejemplobox}[Ejemplo Resuelto 3: Generador Infinitesimal de Traslaciones]
    \textbf{Enunciado:} El operador de traslación espacial finita por una distancia $a \in \mathbb{R}$ se define como:
    \begin{equation}
        \hat{T}(a) \equiv \exp\left(-\frac{\iu}{\hbar}a\hat{p}_x\right).
    \end{equation}
    \begin{enumerate}
        \item Utilice el Lema de Hadamard para demostrar que el operador posición transformado es $\hat{T}\adj(a)\hat{x}\hat{T}(a) = \hat{x} + a\hat{I}$.
        \item Aplique la definición en serie de potencias de $\hat{T}(a)$ sobre una función de onda arbitraria diferenciable $\psi(x)$ en la representación de coordenadas y demuestre que genera una traslación rígida: $\hat{T}(a)\psi(x) = \psi(x - a)$.
    \end{enumerate}

    \textbf{Solución Paso a Paso:}

    \textit{Paso 1: Aplicación del Lema de Hadamard:}
    Como $\hat{p}_x$ es hermítico, $\hat{T}\adj(a) = \exp\left(\frac{\iu}{\hbar}a\hat{p}_x\right)$. Queremos calcular:
    \begin{equation}
        \hat{x}' \equiv \hat{T}\adj(a)\hat{x}\hat{T}(a) = \exp\left(\frac{\iu}{\hbar}a\hat{p}_x\right)\hat{x}\exp\left(-\frac{\iu}{\hbar}a\hat{p}_x\right).
    \end{equation}
    Identificamos en la fórmula de Hadamard \eqref{eq:hadamard_lemma}: el operador conjugante es $\hat{A} = \frac{\iu}{\hbar}a\hat{p}_x$, el operador a transformar es $\hat{B} = \hat{x}$, y el parámetro auxiliar es $\lambda = 1$.
    Calculamos el primer conmutador:
    \begin{equation}
        [\hat{A}, \hat{B}] = \left[\frac{\iu}{\hbar}a\hat{p}_x, \, \hat{x}\right] = \frac{\iu a}{\hbar}[\hat{p}_x, \hat{x}].
    \end{equation}
    Recordando que el conmutador fundamental es $[\hat{x}, \hat{p}_x] = \iu\hbar\hat{I}$, su opuesto es $[\hat{p}_x, \hat{x}] = -\iu\hbar\hat{I}$. Por lo tanto:
    \begin{equation}
        [\hat{A}, \hat{B}] = \frac{\iu a}{\hbar}(-\iu\hbar\hat{I}) = -\iu^2 a \hat{I} = a\hat{I}.
    \end{equation}
    Calculamos el segundo conmutador de la serie de Hadamard:
    \begin{equation}
        [\hat{A}, [\hat{A}, \hat{B}]] = \left[\frac{\iu}{\hbar}a\hat{p}_x, \, a\hat{I}\right] = 0,
    \end{equation}
    puesto que $a\hat{I}$ es un múltiplo escalar de la identidad y conmuta con cualquier operador. Por consiguiente, \textbf{todos los conmutadores de orden dos y superior se anulan}. Sustituyendo en Hadamard:
    \begin{equation}
        \hat{T}\adj(a)\hat{x}\hat{T}(a) = \hat{x} + [\hat{A}, \hat{B}] + 0 = \hat{x} + a\hat{I}.
    \end{equation}
    Esto demuestra de manera puramente operacional que el operador momento $\hat{p}_x$ es el generador infinitesimal de traslaciones espaciales en la posición.

    \textit{Paso 2: Acción sobre la Función de Onda:}
    En la representación de posiciones, el operador momento actúa como la derivada espacial $\hat{p}_x \to -\iu\hbar\frac{\dif}{\dif x}$. Sustituyendo esto en el operador exponencial:
    \begin{equation}
        \hat{T}(a) = \exp\left(-\frac{\iu a}{\hbar}\left(-\iu\hbar\frac{\dif}{\dif x}\right)\right) = \exp\left(\iu^2 a\frac{\dif}{\dif x}\right) = \exp\left(-a\frac{\dif}{\dif x}\right).
    \end{equation}
    Expandiendo en su serie de Taylor operacional:
    \begin{equation}
        \hat{T}(a)\psi(x) = \sum_{n=0}^\infty \frac{1}{n!}\left(-a\frac{\dif}{\dif x}\right)^n \psi(x) = \sum_{n=0}^\infty \frac{(-a)^n}{n!}\frac{\dif^n\psi(x)}{\dif x^n}.
    \end{equation}
    Esta suma infinita es precisamente la \textbf{serie de Taylor} de la función $\psi(x)$ desplazada en una distancia $-a$:
    \begin{equation}
        \hat{T}(a)\psi(x) = \psi(x - a)\footnote{Esta expansión en serie de Taylor asume rigurosamente que la función de onda $\psi(x)$ es analítica (por ejemplo, funciones gaussianas multiplicadas por polinomios), garantizando que su serie de potencias converge a $\psi(x-a)$. Por el Teorema de Stone, el operador traslación $\hat{T}(a) = \eu^{-\iu a\hat{p}_x/\hbar}$ queda unívocamente definido como unitario sobre la totalidad de $\mathcal{L}^2(\mathbb{R})$, aplicándose válidamente incluso sobre funciones no analíticas.}.
    \end{equation}
    ¡El resultado conecta de manera deslumbrante el cálculo diferencial clásico con el álgebra de operadores cuánticos!
\end{ejemplobox}

\begin{figure}[H]
\centering
\begin{subfigure}[t]{0.47\textwidth}
\centering
\begin{tikzpicture}
\begin{axis}[width=\linewidth,height=5.8cm,xlabel={$x$},
  xmin=-4,xmax=6,ymin=0,ymax=1.3,
  legend style={font=\scriptsize,at={(0.5,-0.28)},anchor=north,legend columns=2},ytick=\empty,axis lines=left]
  \addplot[uniblue,very thick,domain=-4:6,samples=200] {exp(-x^2/2)};
  \addlegendentry{$\psi(x)$}
  \addplot[unired,very thick,domain=-4:6,samples=200] {exp(-(x-2)^2/2)};
  \addlegendentry{$\hat T(a)\psi=\psi(x-a)$}
  \draw[-{Stealth},thick] (axis cs:0,1.08) -- (axis cs:2,1.08)
        node[midway,above] {$a>0$};
\end{axis}
\end{tikzpicture}
\caption{Con $a=2$: la funci\'on se desplaza hacia la \textbf{derecha}.}
\end{subfigure}\hfill
\begin{subfigure}[t]{0.52\textwidth}
\centering
\begin{tikzpicture}
\begin{axis}[width=\linewidth,height=5.8cm,xlabel={$x$},
  xmin=-3,xmax=4,ymin=-0.9,ymax=1.5,
  legend style={font=\scriptsize,at={(0.5,-0.28)},anchor=north,legend columns=3},
  grid=both,minor tick num=1]
  \addplot[black,very thick,domain=-3:4,samples=200] {exp(-(x-1)^2/2)};
  \addlegendentry{exacta $\psi(x-1)$}
  \addplot[uniblue,dashed,thick,domain=-3:4,samples=200] {exp(-x^2/2)};
  \addlegendentry{$N=0$}
  \addplot[unired,thick,domain=-3:4,samples=200] {exp(-x^2/2)*(1+x)};
  \addlegendentry{$N=1$}
  \addplot[unigreen,thick,domain=-3:4,samples=200]
    {exp(-x^2/2)*(1+x+(x^2-1)/2)};
  \addlegendentry{$N=2$}
  \addplot[orange,thick,domain=-3:4,samples=200]
    {exp(-x^2/2)*(1+x+(x^2-1)/2+(x^3-3*x)/6)};
  \addlegendentry{$N=3$}
\end{axis}
\end{tikzpicture}
\caption{Sumas parciales de $\sum_{n}\frac{(-a)^n}{n!}\psi^{(n)}(x)$, con $a=1$ y $\psi=\eu^{-x^2/2}$.}
\end{subfigure}
\caption{Acci\'on de $\hat T(a)=\eu^{-a\,\dif/\dif x}$: el desarrollo de Taylor reconstruye $\psi(x-a)$.}
\label{fig:traslacion}
\end{figure}

\subsection{Simetrías Continuas, Leyes de Conservación y el Teorema de Noether Cuántico}
\label{subsec:noether_cuantico}

El concepto de operador exponencial desarrollado en esta sección constituye la piedra angular de la física teórica moderna: \textbf{conecta directamente las simetrías geométricas del espacio-tiempo con las leyes de conservación de magnitudes físicas} \cite{zettili, sakurai}.

\subsubsection{Transformaciones Unitarias Infinitesimales}
Consideremos una transformación continua parametrizada por un escalar infinitesimal $\epsilon \ll 1$. Dicha transformación se implementa mediante un operador unitario $\hat{U}_\epsilon$ cercano a la identidad:
\begin{equation}
    \hat{U}_\epsilon = \hat{I} - \frac{\iu}{\hbar}\epsilon\,\hat{G},
    \label{eq:infinitesimal_unitary}
\end{equation}
donde $\hat{G}$ es un operador \textbf{hermítico} ($\hat{G}\adj = \hat{G}$), denominado el \textbf{generador de la transformación}. El factor $-\iu/\hbar$ se introduce por convención para que $\hat{G}$ posea las dimensiones físicas de la magnitud observable correspondiente y para garantizar la unitariedad:
\begin{equation}
    \hat{U}_\epsilon\adj \hat{U}_\epsilon = \left(\hat{I} + \frac{\iu}{\hbar}\epsilon\,\hat{G}\right)\left(\hat{I} - \frac{\iu}{\hbar}\epsilon\,\hat{G}\right) = \hat{I} + \mathcal{O}(\epsilon^2).
\end{equation}

Bajo esta transformación, un estado $\ket{\psi}$ y un observable $\hat{A}$ se modifican infinitesimalmente de acuerdo a:
\begin{align}
    \ket{\psi'} &= \hat{U}_\epsilon \ket{\psi} = \ket{\psi} - \frac{\iu}{\hbar}\epsilon\,\hat{G}\ket{\psi}, \\
    \hat{A}' &= \hat{U}_\epsilon\adj \hat{A} \hat{U}_\epsilon = \left(\hat{I} + \frac{\iu}{\hbar}\epsilon\,\hat{G}\right)\hat{A}\left(\hat{I} - \frac{\iu}{\hbar}\epsilon\,\hat{G}\right) \nonumber\\
    &= \hat{A} + \frac{\iu}{\hbar}\epsilon\left(\hat{G}\hat{A} - \hat{A}\hat{G}\right) + \mathcal{O}(\epsilon^2) = \hat{A} + \frac{\iu}{\hbar}\epsilon\,[\hat{G}, \hat{A}] + \mathcal{O}(\epsilon^2).
\end{align}
La variación infinitesimal del observable $\delta\hat{A} \equiv \hat{A}' - \hat{A}$ viene dada estrictamente por:
\begin{equation}
    \boxed{\delta\hat{A} = +\frac{\iu}{\hbar}\epsilon\,[\hat{G}, \hat{A}]}.
    \label{eq:infinitesimal_variation}
\end{equation}
(Nótese que para una traslación espacial con $\hat{G} = \hat{p}_x$ y $\epsilon = a$, tenemos $[\hat{p}_x, \hat{x}] = -\iu\hbar\hat{I}$, lo cual conduce de forma plenamente consistente a $\delta\hat{x} = +\frac{\iu}{\hbar}a(-\iu\hbar\hat{I}) = +a\hat{I}$, recuperando $\hat{x}' = \hat{x} + a\hat{I}$).

\subsubsection{Imagen de Heisenberg, Invariancia del Hamiltoniano y Constantes de Movimiento}
Supongamos que la dinámica de un sistema aislado está gobernada por un Hamiltoniano $\hat{H}$ independiente del tiempo. Existen dos formas equivalentes de describir la evolución temporal cuántica:
\begin{itemize}
    \item \textbf{Imagen de Schrödinger:} Los estados evolucionan en el tiempo mediante el operador unitario $\hat{U}(t) = \eu^{-\iu\hat{H}t/\hbar}$ ($\ket{\psi_S(t)} = \hat{U}(t)\ket{\psi(0)}$), mientras que los operadores observables $\hat{A}_S$ permanecen fijos.
    \item \textbf{Imagen de Heisenberg:} Los estados permanecen fijos ($\ket{\psi_H} = \ket{\psi(0)}$) y la dependencia temporal se traslada a los operadores observables mediante la transformación de semejanza unitaria:
    \begin{equation}
        \hat{A}_H(t) \equiv \hat{U}\adj(t)\hat{A}_S\hat{U}(t) = \eu^{+\frac{\iu}{\hbar}\hat{H}t}\,\hat{A}_S\,\eu^{-\frac{\iu}{\hbar}\hat{H}t},
    \end{equation}
    de modo que el valor esperado físico es idéntico en ambas imágenes: $\matrixel{\psi_S(t)}{\hat{A}_S}{\psi_S(t)} = \matrixel{\psi(0)}{\hat{A}_H(t)}{\psi(0)} = \ev{\hat{A}(t)}$.
\end{itemize}
Observemos que $\hat{A}_H(t) = \eu^{+\frac{\iu}{\hbar}\hat{H}t}\hat{A}_S\eu^{-\frac{\iu}{\hbar}\hat{H}t}$ posee exactamente la estructura del \textbf{Lema de Hadamard} \eqref{eq:hadamard_ansatz} con $\lambda = t$ y $\hat{A} \to \frac{\iu}{\hbar}\hat{H}$. Derivando respecto al tiempo $t$ mediante la Ec.~\eqref{eq:hadamard_first_diff} (e incluyendo cualquier dependencia temporal explícita $\partial\hat{A}_S/\partial t$):
\begin{equation}
    \frac{\dif\hat{A}_H(t)}{\dif t} = \left[\frac{\iu}{\hbar}\hat{H}, \, \hat{A}_H(t)\right] + \left(\frac{\partial\hat{A}_S}{\partial t}\right)_H \implies \boxed{\frac{\dif\hat{A}_H}{\dif t} = \frac{1}{\iu\hbar}[\hat{A}_H, \hat{H}] + \left(\frac{\partial\hat{A}}{\partial t}\right)_H}.
    \label{eq:heisenberg_eq_motion}
\end{equation}
Esta es la célebre \textbf{Ecuación de Movimiento de Heisenberg} (que aplicaremos en el Problema Resuelto 4 para deducir la precesión de Larmor).

Se dice que una transformación unitaria $\hat{U}$ generada por $\hat{G}$ representa una \textbf{simetría del sistema} si el Hamiltoniano es invariante ante la transformación:
\begin{equation}
    \hat{H}' = \hat{U}\adj \hat{H} \hat{U} = \hat{H} \iff [\hat{G}, \hat{H}] = 0.
    \label{eq:hamiltonian_invariance}
\end{equation}
Tomando el valor esperado de la Ec.~\eqref{eq:heisenberg_eq_motion} para el generador $\hat{G}$ en un estado arbitrario $\ket{\psi}$, obtenemos el \textbf{Teorema de Ehrenfest}:
\begin{equation}
    \frac{\dif}{\dif t}\ev{\hat{G}} = \frac{1}{\iu\hbar}\ev{[\hat{G}, \hat{H}]} + \left\langle \frac{\partial\hat{G}}{\partial t} \right\rangle.
    \label{eq:ehrenfest_theorem}
\end{equation}
Si el generador $\hat{G}$ no depende explícitamente del tiempo ($\partial\hat{G}/\partial t = 0$) y conmuta con el Hamiltoniano ($[\hat{G}, \hat{H}] = 0$), se deduce de forma inmediata:
\begin{equation}
    \frac{\dif}{\dif t}\ev{\hat{G}} = 0 \implies \boxed{\ev{\hat{G}} = \text{Constante en el tiempo}}.
\end{equation}

\begin{teoremabox}[Teorema de Noether Cuántico]
    A toda simetría continua del Hamiltoniano le corresponde un observable físico hermítico cuya media estadística es estrictamente constante en el tiempo (una \textbf{constante de movimiento}).
\end{teoremabox}

\subsubsection{El Puente Clásico-Cuántico: De los Corchetes de Poisson a los Conmutadores}
En la formulación hamiltoniana de la mecánica clásica, el estado de una partícula se describe en el \textbf{espacio de fase} de coordenadas y momentos $(x_i, p_i)$ (formalizado geométricamente como una variedad simpléctica dotada de la 2-forma $\omega = \dif x_i \wedge \dif p_i$; no es necesario conocer geometría diferencial para seguir este curso, basta con recordar las derivadas parciales ordinarias respecto a $x_i$ y $p_i$). La evolución temporal de una magnitud clásica $A(x, p)$ está gobernada por los \textbf{Corchetes de Poisson}:
\begin{equation}
    \frac{\dif A}{\dif t} = \{A, H\}_{\text{clásico}} + \frac{\partial A}{\partial t}, \qquad \{A, B\}_{\text{clásico}} \equiv \sum_{i=1}^3 \left( \frac{\partial A}{\partial x_i}\frac{\partial B}{\partial p_i} - \frac{\partial A}{\partial p_i}\frac{\partial B}{\partial x_i} \right).
\end{equation}
El célebre \textbf{Postulado de Cuantización Canónica de Dirac} establece que la estructura geométrica simpléctica clásica se traduce en la estructura algebraica de conmutadores en el espacio de Hilbert mediante el mapeo universal:
\begin{equation}
    \boxed{\{A, B\}_{\text{clásico}} \quad \xrightarrow{\quad \text{Cuantización Canónica} \quad} \quad \frac{1}{\iu\hbar}[\hat{A}, \hat{B}]}.
    \label{eq:dirac_canonical_quantization}
\end{equation}
Toda magnitud conservada en la mecánica clásica ($\{G, H\} = 0$) se transforma en cuántica en un observable compatible con la energía ($[\hat{G}, \hat{H}] = 0$), tal como se ilustra sintéticamente en la Figura~\ref{fig:simetrias}.

\begin{figure}[H]
\centering
\begin{tikzpicture}[x=1cm,y=1cm]
  \node[font=\small\bfseries] at (1.0,0.9) {Simetr\'ia};
  \node[font=\small\bfseries] at (5.2,0.9) {Generador};
  \node[font=\small\bfseries] at (9.0,0.9) {Magnitud conservada};
  \draw[uniblue,thick] (-0.8,0.55) -- (11,0.55);
  \begin{scope}[shift={(0,-0.6)}]
    \draw[thick,uniblue] (0,0) circle (0.45);
    \draw[thick,uniblue] (0,0) -- (0,0.3);
    \draw[thick,uniblue] (0,0) -- (0.22,-0.1);
    \draw[-{Stealth},thick,unired] (0.7,0.1) -- (1.2,0.1) node[right,font=\small,text=black,xshift=2pt] {$t\to t+\Delta t$};
    \node at (5.2,0) {$\hat H$};
    \node at (9.0,0) {Energ\'ia $E$};
  \end{scope}
  \draw[gray!50] (-0.8,-1.3) -- (11,-1.3);
  \begin{scope}[shift={(0,-2.1)}]
    \draw[thick,uniblue,fill=uniblue!15] (-0.3,-0.2) rectangle (0.1,0.2);
    \draw[thick,uniblue,dashed] (0.5,-0.2) rectangle (0.9,0.2);
    \draw[-{Stealth},thick,unired] (0.15,0.38) -- (0.85,0.38);
    \node[font=\small,right] at (1.3,0) {$\vec r\to\vec r+\vec a$};
    \node at (5.2,0) {$\hat{\vec p}$};
    \node at (9.0,0) {Momento lineal $\vec p$};
  \end{scope}
  \draw[gray!50] (-0.8,-2.8) -- (11,-2.8);
  \begin{scope}[shift={(0,-3.6)}]
    \draw[thick,uniblue] (0,0) circle (0.4);
    \fill[uniblue] (0.4,0) circle (0.05);
    \draw[-{Stealth},thick,unired] (0.55,0.1) arc (10:110:0.55);
    \node[font=\small,right] at (1.3,0) {$\theta\to\theta+\Delta\theta$};
    \node at (5.2,0) {$\hat{\vec L}$};
    \node at (9.0,0) {Momento angular $\vec L$};
  \end{scope}
  \draw[uniblue,thick] (-0.8,-4.3) -- (11,-4.3);
  \foreach \y in {-0.6,-2.1,-3.6}{
    \draw[-{Stealth},gray] (3.3,\y) -- (4.3,\y);
    \draw[-{Stealth},gray] (6.2,\y) -- (7.0,\y);
  }
  \node[font=\scriptsize,unigray,align=center] at (5.1,-4.8)
    {Homogeneidad del tiempo $\;\cdot\;$ homogeneidad del espacio $\;\cdot\;$ isotrop\'ia del espacio};
\end{tikzpicture}
\caption{Simetr\'ias continuas, generadores herm\'iticos y magnitudes conservadas (Teorema de Noether cu\'antico).}
\label{fig:simetrias}
\end{figure}
